% ============================================================
% SECTION: High-Fidelity Tactical Logic Validation (HTLV)
% ============================================================

\subsection{High-Fidelity Tactical Logic Validation (HTLV)}
\label{sec:htlv_architecture}

Following the \gls{msla} conceptual experiments, this section presents the \gls{htlv}: a
prototype high-fidelity environment incorporating NASA's DAIDALUS as a modeled \gls{rwc}
reasoning engine. The \gls{htlv} evaluates kinematic alignment with RTCA DO-365B~\cite{rtca_do365b}
and ASTM F3442~\cite{astm_f3442} separation guidelines through precision kinematic updates at
10\,Hz (physics step), with alert-level evaluation running continuously over a simulated 1-hour
mission window. As a research initiative, it serves to investigate rather than certify complete
operational compliance.

\textbf{Methodological Note:} To ensure the precise sub-second physics iteration required for
DAIDALUS validation, this tactical micro-simulator was custom-developed in a
high-performance compiled language. It remained in active research development throughout
Point D and is designed to become a candidate high-performance UTM/\gls{xtm} research tool
for ITA and \gls{icea}.

% -----------------------------------------------------------
\subsubsection{HTLV Architecture and Execution Engine}
% -----------------------------------------------------------

The \gls{htlv} operates as a unified simulation runner coupling standard geometric avoidance
schemes with advanced \gls{rwc} algorithms:

\begin{itemize}
    \item \textbf{Simulator Core:} Built on a high-performance deterministic engine, scenarios
    are generated from structured configuration files. A custom 3-D spatial grid algorithm
    accelerates neighbor search and evaluation of the 20\,m horizontal / 10\,m vertical
    MACproxy separation threshold.

    \item \textbf{DAIDALUS Bridge:} The NASA \gls{daa} logic is compiled as a native module
    directly integrated into the simulation pipeline. The engine continuously queries DAIDALUS
    for relative aircraft pairings, extracting alert levels, time-to-violation (TTV), and
    horizontal safe-heading resolution bands.

    \item \textbf{Monitoring vs.\ Reactive Steering:} The setup physically decouples the
    \emph{DAA Monitoring System} (which populates active conflict logs) from the
    \emph{Reactive Avoidance System} (which translates DAIDALUS guidance bands into physical
    control inputs). This separation enables distinct validation of standard alerts
    (Preventive, Corrective, Warning) independent of the aircraft's kinematic evasive
    response.
\end{itemize}

% -----------------------------------------------------------
\subsubsection{Methodological Comparison: Macro-Scale vs.\ High-Fidelity}
\label{subsec:rust_vs_python}
% -----------------------------------------------------------

Table~\ref{tab:rust_comparison} highlights the complementary strengths of the two modelling
approaches used throughout the experimental campaign.

\begin{table}[H]
\centering
\caption{Macro-Scale Simulator vs.\ High-Fidelity Tactical Engine}
\label{tab:rust_comparison}
\renewcommand{\arraystretch}{1.2}
\resizebox{\textwidth}{!}{%
\begin{tabular}{lp{6.5cm}p{6.5cm}}
\toprule
\textbf{Dimension} & \textbf{Macro-Scale (MSLA)} & \textbf{High-Fidelity Engine (HTLV)} \\
\midrule
Physics Engine    & Simple kinematics (dt = 1\,s)            & High-performance deterministic engine (10\,Hz) \\
\gls{daa} Engine  & Custom heuristic (proximity-based)        & NASA DAIDALUS (DO-365B aligned) \\
Alert Information & Binary (Threat / No Threat)               & Multi-level (Preventive / Corrective / Warning), TTV, Heading Bands \\
Avoidance Maneuver & Fixed angular (60--90°) heuristic        & DAIDALUS Bands + Horiz./Vert.\ offsets \\
Scale             & Up to 2\,000 concurrent UAS               & Up to 1\,000 concurrent UAS (1-hour window) \\
Simulation Speed  & Near-realtime (up to 1\,600$\times$)      & Dependent on conflict graph size (1--6\,500$\times$) \\
\bottomrule
\end{tabular}%
}
\end{table}

\begin{figure}[htbp]
    \centering
    \includegraphics[width=1.0\textwidth]{Docs_ref/Rust_sim/Saturated_500_SJC.png}
    \caption{Visualization of a saturated airspace simulation around SJC Airport with $N=500$
    simultaneous drones operating under DAIDALUS guidance. The trajectory entanglement visible
    in the dense corridor illustrates the combinatorial interaction growth that drives the
    scalability limits detailed in Section~\ref{sec:scen2_sweep}.}
    \label{fig:rust_htlv_viewer}
\end{figure}

% -----------------------------------------------------------
\subsubsection{Complementary Strengths and Integration Path}
\label{subsec:strengths_integration}
% -----------------------------------------------------------

The \textbf{MACproxy Macro-layer} effectively assesses the three-layer conflict management
model defined in the \gls{conops} --- Strategic (\gls{xtm}), Tactical (\gls{daa}), and Collision
Avoidance. It utilises real SJC geography and produces metrics directly applicable to capacity
scoping with large drone populations.

The \textbf{HTLV Micro-layer} provides the increased fidelity level necessary to investigate
algorithmic alignment with international guidelines such as RTCA DO-365B~\cite{rtca_do365b}
and ASTM F3442~\cite{astm_f3442}. Its localized physics engine resolves finer kinematic details
during simulated \gls{rwc} avoidance maneuvers.

% -----------------------------------------------------------
\subsubsection{DAIDALUS Core Logic}
\label{sec:daidalus}
% -----------------------------------------------------------

The DAIDALUS algorithm constitutes the mathematical backbone of the \gls{htlv}. Within the
\gls{daa} layered architecture, DAIDALUS implements the \gls{rwc} function as defined in
RTCA DO-365B~\cite{rtca_do365b}: it transforms surveillance data into maneuver guidance
(\emph{Bands}) for the \gls{rpic} or autonomous system, with the objective of maintaining the
Well Clear separation volume.

\textbf{Justification for Evaluation:} DAIDALUS was selected for integration because it
serves as the globally recognised aviation-authority reference implementation for DO-365B.
Evaluating it inside the SIMUA parameters answers a critical regulatory question:
\emph{Can standardised mathematical compliance models developed for sparse airspace scale safely
into Brazilian high-density urban corridors?} Importantly, DAIDALUS operates strictly on
immediate kinematic state vectors (current position and velocity) rather than tracking
pre-planned mission routes. This fundamental characteristic decouples tactical collision
avoidance from strategic 4-D paths, making it the ideal testbed to investigate how onboard
autonomy impacts overarching fleet operations.

\paragraph{Cooperative vs.\ Non-Cooperative Operational Modes}
In a cooperative airspace, DAIDALUS operates with high-fidelity telemetry data (\gls{ads-b}),
assessing the \emph{Well-Clear} timeline using exact time-to-convergence inequalities based
on minimum vertical/horizontal separation thresholds~\cite{wellclear_nasa}. The strategic
differentiator of the Point D architecture is its planned behaviour when fed by onboard sensors
(Radar and Vision) detecting non-cooperative aircraft, which provide data characterised by
noise and systematic measurement errors.

\paragraph{Planned Extension: Sensor Uncertainty Mitigation (SUM) and VISA Smoothing}
\label{sec:sum_visa_future}
The DAIDALUS framework natively supports \gls{sum}
(\emph{Sensor Uncertainty Mitigation})~\cite{applying_sum_nasa}, which uses error
distributions ($z$-scores) to dynamically expand the Well Clear volume into a safety
ellipsoid encompassing the intruder's ``worst probable position,'' and \gls{visa}
(\emph{Virtual Intruder State Aggregation}), which smooths temporal jitter by averaging
consecutive sensor updates~\cite{sum_evaluation_nasa}.

\textbf{Current experimental scope:} All scalability runs presented in
Section~\ref{sec:scen2_sweep} operate under \emph{cooperative perfect-state} assumptions
(exact position and velocity vectors), so neither SUM expansion nor VISA smoothing was
exercised in the results reported here. Activating SUM with realistic onboard-sensor noise
profiles (mmWave Radar, Computer Vision) is a programmed next step that will widen the
effective Well Clear volume and is expected to increase both alert rates and route
inefficiency --- quantifying that penalty is a priority for the next project phase.

\paragraph{Planned Extension: Human-in-the-Loop (HITL) Alerting Validation}
The DAIDALUS alert taxonomy maps to absolute reaction timelines relevant for \gls{rpic}
authority~\cite{atc_acceptability_nasa}: \textbf{Preventive} ($\sim$60\,s to \gls{cpa}),
\textbf{Corrective} ($\sim$30\,s), and \textbf{Warning} ($\sim$15\,s). While these thresholds
are configured in the current \gls{htlv} engine, the scalability sweep does not yet inject
\gls{rpic} reaction latency or decision-loop delays into the simulation. Validating that the
5--15\,s pilot network latency budget~\cite{atc_acceptability_nasa} is safely absorbed under
dense traffic conditions remains an explicit future-work objective.

\paragraph{DAIDALUS Steering Strategies Under Evaluation}
\label{sec:action_modes}
Two complementary DAIDALUS-mediated steering strategies were evaluated in the scalability sweep:
\begin{enumerate}
    \item \textbf{Continuous Safe Band:} At every physics step, the engine selects the nearest
    heading within the current DAIDALUS \emph{safe band} (i.e.\ the set of headings that
    do not violate the Well Clear volume for any tracked intruder) and blends that heading
    continuously into the active trajectory. This approach produces the smoothest flight path
    and maximises compliance with the full DAIDALUS resolution logic.

    \item \textbf{Preferred Horizontal Resolution:} Rather than blending across the full safe
    band, the engine hard-couples to the single optimal avoidance heading recommended by the
    DAIDALUS \emph{preferred resolution} output at each update. This yields a more decisive
    turn toward the algorithmically chosen escape direction.
\end{enumerate}

\noindent Both modes are compared against a \textbf{No-DAIDALUS baseline} that uses a simple
proximity-triggered angular maneuver (the same heuristic used in the \gls{msla}), providing
the reference ``blind flight'' degradation curve.

% -----------------------------------------------------------
\subsubsection{Scenario 2: DAIDALUS Scalability Evaluation}
\label{sec:scen2_sweep}
% -----------------------------------------------------------

To push beyond single-encounter validation, a large-scale tactical resolution sweep was
executed within the \gls{htlv} unified runner. This benchmark compares the no-DAIDALUS
baseline against the two DAIDALUS steering strategies under fixed geographic and temporal
parameters over a 1-hour simulation window. Fleet densities ranged from $N = 50$ to
$N = 1{,}000$ simultaneous drones in the geofenced corridor.

\paragraph{Full Cross-Density Results}

Table~\ref{tab:scen2_full} presents the complete sweep across all density levels. The metrics
are: \textbf{CR} --- mission completion ratio (completed\,/\,scheduled);
\textbf{MAC} --- total MACproxy events (20\,m horizontal / 10\,m vertical violations);
\textbf{Ineff.} --- route inefficiency percentage (excess distance relative to ideal chord);
\textbf{Wall} --- simulator wall-clock time.

\begin{table}[htbp]
\caption{DAIDALUS Scalability Sweep --- Scenario 2, 1-hour window, 10\,Hz physics,
$N = 50$ to $N = 1{,}000$.}
\label{tab:scen2_full}
\centering
\renewcommand{\arraystretch}{1.2}
\resizebox{\textwidth}{!}{%
\begin{tabular}{llrrrr}
\toprule
\textbf{Mode} & $N$ & \textbf{CR} & \textbf{MAC} & \textbf{Ineff.\ (\%)} & \textbf{Wall (s)} \\
\midrule
\textbf{Preferred Horizontal Resolution}
 & 50   & 0.853 & 0      & 6.22   & 4.6 \\
 & 100  & 0.856 & 0      & 8.33   & 14.5 \\
 & 200  & 0.851 & 2      & 11.85  & 49.9 \\
 & 500  & 0.774 & 36     & 49.73  & 623.3 \\
 & 1000 & 0.306 & 1\,202  & 90.85  & 11\,447.9 \\
\addlinespace
\textbf{Continuous Safe Band}
 & 50   & 0.853 & 0      & 3.68   & 3.0 \\
 & 100  & 0.860 & 5      & 6.61   & 12.7 \\
 & 200  & 0.857 & 33     & 14.98  & 65.7 \\
 & 500  & 0.724 & 1\,619  & 66.46  & 1\,574.6 \\
 & 1000 & 0.334 & 19\,050 & 138.48 & 14\,721.0 \\
\addlinespace
\textbf{Baseline (No DAIDALUS)}
 & 50   & 0.849 & 0      & 4.41   & 0.5 \\
 & 100  & 0.856 & 0      & 8.16   & 0.9 \\
 & 200  & 0.843 & 0      & 19.74  & 2.3 \\
 & 500  & 0.774 & 42     & 70.17  & 13.2 \\
 & 1000 & 0.460 & 3\,601  & 145.56 & 155.5 \\
\bottomrule
\end{tabular}%
}
\end{table}


% -----------------------------------------------------------
\paragraph{Head-to-Head: DAIDALUS vs.\ No-DAIDALUS Baseline}
\label{sec:daidalus_vs_baseline}
% -----------------------------------------------------------

Table~\ref{tab:daidalus_vs_baseline} distills the direct comparison between the best DAIDALUS
mode at each density level (Preferred Horizontal Resolution) and the proximity-only baseline
to highlight where standard-aligned RWC guidance adds value --- and where it becomes
counterproductive.

\begin{table}[htbp]
\caption{DAIDALUS (Preferred Horizontal Resolution) vs.\ No-DAIDALUS Baseline --- direct
  comparison at each density level.  $\Delta$Ineff.\ shows the route inefficiency
  \emph{reduction} provided by DAIDALUS (positive = DAIDALUS is more efficient).}
\label{tab:daidalus_vs_baseline}
\centering
\renewcommand{\arraystretch}{1.2}
\resizebox{\textwidth}{!}{%
\begin{tabular}{r|rr|rr|rr}
\toprule
 & \multicolumn{2}{c|}{\textbf{MAC Events}}
 & \multicolumn{2}{c|}{\textbf{Route Ineff.\ (\%)}}
 & \multicolumn{2}{c}{\textbf{Completion Ratio}} \\
$N$ & \textbf{Baseline} & \textbf{DAIDALUS}
    & \textbf{Baseline} & \textbf{DAIDALUS}
    & \textbf{Baseline} & \textbf{DAIDALUS} \\
\midrule
50   & 0     & 0     & 4.41  & 6.22  & 0.849 & 0.853 \\
100  & 0     & 0     & 8.16  & 8.33  & 0.856 & 0.856 \\
200  & 0     & 2     & 19.74 & 11.85 & 0.843 & 0.851 \\
500  & 42    & 36    & 70.17 & 49.73 & 0.774 & 0.774 \\
1000 & 3\,601 & 1\,202 & 145.56 & 90.85 & 0.460 & 0.306 \\
\bottomrule
\end{tabular}%
}
\end{table}

Three distinct operational regimes emerge from this comparison:

\textbf{(i) Low density ($N \leq 100$):\ equivalent safety, comparable efficiency.}
Both modes record zero MACproxy events, similar completion rates ($\approx$85\%), and
nearly identical route inefficiency. At these densities the encounter rate is low enough
that the simple proximity heuristic suffices; DAIDALUS adds algorithmic rigour and
standard-aligned alerting but does not yet differentiate operationally in the metrics
measured.

\textbf{(ii) Moderate density ($N = 200$--$500$):\ DAIDALUS improves route quality.}
This is the regime where DAIDALUS delivers its clearest advantage. At $N = 200$, the
baseline's proximity heuristic generates \textbf{19.7\% route inefficiency} while DAIDALUS
achieves \textbf{11.9\%} --- a $\approx$40\% reduction in excess path distance --- with
comparable safety (2 MACs vs.\ 0). At $N = 500$, DAIDALUS reduces MAC events from 42 to 36
\emph{and} simultaneously cuts route inefficiency from 70\% to 50\%. The advantage arises
because DAIDALUS computes the \emph{minimum-deviation safe heading} via the Well Clear
inequality, whereas the proximity heuristic applies a fixed angular offset that often
over-corrects. For battery-constrained sUAS, this 20--30 percentage-point route saving
translates directly into extended operational range and mission feasibility.

\textbf{(iii) High density ($N = 1{,}000$):\ reactive saturation crossover.}
At extreme density, the comparison inverts. DAIDALUS reduces MAC events (1,202 vs.\ 3,601
for the baseline) but at a severe cost: completion ratio drops to 30.6\% versus the
baseline's 46.0\%, and the Continuous Safe Band variant generates \textbf{19,050} MAC
events --- \emph{5.3$\times$ worse than no guidance at all}. This paradoxical inversion
signals the onset of \emph{reactive oscillation deadlock}.

% -----------------------------------------------------------
\paragraph{Root Cause: Reactive Oscillation Deadlock at Scale}
\label{sec:daidalus_saturation}
% -----------------------------------------------------------

The wall-clock data reveals the physical mechanism behind the high-density collapse.
Continuous Safe Band requires 14,721\,s ($\approx$4~hours) to simulate a 3,600\,s scenario
at $N = 1{,}000$, an amplification factor of $4.1\times$ the simulated time, compared to
$0.04\times$ for the baseline. This growth is consistent with the $O(N^2)$ structure of
the DAIDALUS pairwise evaluation graph: at $N = 1{,}000$, the engine evaluates
\textbf{102,655 alert pairs} per step versus 74,055 for the baseline.

The continuous re-querying triggers \emph{heading oscillations} across multiple interacting
pairs: each drone adjusts heading in response to newly generated conflicts caused by the
previous step's adjustments, creating a feedback loop that prevents mission completion and
paradoxically \emph{increases} the total number of separation violations. Preferred
Horizontal Resolution exhibits qualitatively similar behaviour (wall-clock: 11,448\,s),
confirming that pairwise re-evaluation dominates regardless of how the resolution heading
is selected.

\textbf{Implication:} DAIDALUS, as a pairwise reactive algorithm designed for sparse
cooperative encounters (DO-365B), was not architecturally intended for dense multi-agent
environments where cascading pair interactions dominate. The scalability sweep empirically
demonstrates the density threshold ($N \approx 300$--$500$ in a 5\,km corridor) at which
reactive logic alone becomes insufficient and strategic pre-separation (\gls{xtm}) becomes
operationally mandatory.


% -----------------------------------------------------------
\paragraph{Performance Trends}
% -----------------------------------------------------------

The following charts illustrate the degradation curves as the number of drones ($N$) increases
across the geofenced corridor.

\begin{figure}[htbp]
    \centering
    \begin{minipage}{0.48\textwidth}
        \centering
        \includegraphics[width=\linewidth]{Docs_ref/Rust_sim/report_sections/figures/scen2_daidalus_sweep/fig_scen2_completion_ratio_vs_n.png}
        \caption{Mission completion ratio vs.\ fleet size. Both DAIDALUS modes sustain above
        80\% completion through $N = 200$, then degrade sharply. At $N = 1{,}000$,
        completion drops below 35\%, indicating the reactive saturation limit.}
        \label{fig:scen2_completion}
    \end{minipage}\hfill
    \begin{minipage}{0.48\textwidth}
        \centering
        \includegraphics[width=\linewidth]{Docs_ref/Rust_sim/report_sections/figures/scen2_daidalus_sweep/fig_scen2_macproxy_vs_n.png}
        \caption{MACproxy count (safety violations) vs.\ fleet size. Continuous Safe Band
        generates 19,050 violations at $N = 1{,}000$ --- exceeding the 3,601 of the
        no-DAIDALUS baseline --- confirming the reactive oscillation deadlock.}
        \label{fig:scen2_macproxy}
    \end{minipage}

    \vspace{0.5cm}
    \begin{minipage}{0.48\textwidth}
        \centering
        \includegraphics[width=\linewidth]{Docs_ref/Rust_sim/report_sections/figures/scen2_daidalus_sweep/fig_scen2_route_inefficiency_vs_n.png}
        \caption{Route inefficiency (\%) vs.\ fleet size. Continuous Safe Band reaches 138\%
        excess distance at $N = 1{,}000$, rendering battery-constrained sUAS missions
        infeasible. Even at $N = 500$, the 50--66\% penalty imposes severe operational
        constraints.}
        \label{fig:scen2_ineff}
    \end{minipage}\hfill
    \begin{minipage}{0.48\textwidth}
        \centering
        \includegraphics[width=\linewidth]{Docs_ref/Rust_sim/report_sections/figures/scen2_daidalus_sweep/fig_scen2_wall_seconds_vs_n.png}
        \caption{Simulator wall-clock time vs.\ fleet size. The $O(N^2)$ computational growth
        of the DAIDALUS pairwise evaluation is clearly visible, reaching over 14,000\,s for
        Safe Band at $N = 1{,}000$ --- a 94$\times$ slowdown relative to the baseline.}
        \label{fig:scen2_wall}
    \end{minipage}
\end{figure}


% -----------------------------------------------------------
\subsubsection{Key Finding: The xTM Imperative}
\label{sec:xtm_imperative}
% -----------------------------------------------------------

The sweep results provide the empirical foundation for a central conclusion: \textbf{DAIDALUS
alone is necessary but not sufficient for high-density UTM operations.}

At low-to-moderate densities ($N \leq 200$), the DAIDALUS-based \gls{rwc} layer effectively
resolves conflicts with near-zero safety violations and acceptable route penalties. This
validates DAIDALUS as a viable tactical backbone for initial BVLOS deployments in the density
regimes targeted by SIMUA.

Beyond $N \approx 300$--$500$, however, the pairwise reactive nature of DAIDALUS encounters
a fundamental scalability ceiling: cascading pair interactions generate oscillatory
heading adjustments that degrade both safety and efficiency below the no-guidance baseline.
This is not a flaw in the DAIDALUS algorithm --- it performs precisely as designed for
sparse encounter geometries --- but rather a confirmation that \textbf{high-density
corridors require strategic pre-separation (\gls{xtm})} to reduce the effective encounter
rate before tactical DAIDALUS logic engages.

This finding directly reinforces the layered architecture proposed in the \gls{conops}:
\begin{enumerate}
    \item \textbf{Strategic layer (\gls{xtm}):} Pre-tactical flow management, route
    deconfliction, and temporal separation to keep the effective $N$ within the DAIDALUS
    operational envelope ($N \leq 200$--$300$ for a 5\,km corridor).
    \item \textbf{Tactical layer (DAIDALUS):} Real-time RWC guidance for residual conflicts
    that escape strategic deconfliction, operating in the density regime where it is proven
    effective.
    \item \textbf{Collision avoidance layer:} Last-resort emergency maneuvers for encounters
    that breach the RWC boundary.
\end{enumerate}

\noindent The MSLA results (Section~\ref{sec:msla}), which independently demonstrate
the effectiveness of xTM in reducing encounter rates at fleet scale, are therefore
\emph{directly complementary} to the HTLV findings: together they validate the full
strategic-to-tactical pipeline.

% -----------------------------------------------------------
\subsubsection{Future Integration Directions}
% -----------------------------------------------------------

The \gls{htlv} is designed as a modular research platform. The results above suggest two
concrete integration directions for subsequent project phases.

First, the DAIDALUS resolution logic should be integrated as the tactical avoidance engine
within the \gls{msla} fleet simulator, replacing the current macroscopic heuristic. This
closed-loop integration will allow quantification of how xTM pre-separation + DAIDALUS
tactical guidance perform together across the full density range, and will empirically
determine the minimum xTM separation requirements that keep DAIDALUS within its effective
operating envelope.

Second, once the multi-sensor benchmark (Section~\ref{sec:multisensor-benchmark}) produces
calibrated detection probability profiles per sensor and lighting condition, those empirical
envelopes can be injected into the HTLV pipeline as the source of intruder state estimates
--- replacing the current idealised perfect-state assumption and enabling end-to-end
evaluation of the full sense-to-avoid chain under realistic uncertainty budgets.

% -----------------------------------------------------------
\subsubsection{Summary for DECEA Evaluation}
\label{subsec:decea_summary}
% -----------------------------------------------------------

The sequential value chain --- from mathematical pre-tactical capacity bounding in SCOPAS,
through macroscopic 4-D strategic reservations (\gls{msla}), down to high-fidelity,
algorithmically aligned tactical deviation (\gls{htlv}) --- models a holistic conceptual
approach for managing \gls{uas} traffic in Brazilian airspace.

The DAIDALUS scalability sweep provides three operationally relevant conclusions for DECEA
assessment:

\begin{enumerate}
    \item \textbf{DAIDALUS is effective as a tactical RWC layer at operationally relevant
    densities.} At $N \leq 200$ (the regime corresponding to initial BVLOS deployments),
    both steering modes achieve near-zero safety violations with completion rates above 85\%
    and route inefficiency below 15\%.

    \item \textbf{A hard scalability ceiling exists for reactive-only guidance.}
    Beyond $N \approx 300$--$500$, DAIDALUS-only control collapses due to cascading pairwise
    oscillations. At $N = 1{,}000$, the system generates more separation violations than the
    no-guidance baseline, confirming that reactive logic alone cannot sustain high-density
    operations.

    \item \textbf{Strategic pre-separation (xTM) is mandatory above the tactical ceiling.}
    The results empirically validate the layered xTM + DAIDALUS architecture proposed in the
    \gls{conops}: xTM keeps the effective encounter rate within the DAIDALUS operational
    envelope, while DAIDALUS handles residual conflicts with proven effectiveness.
\end{enumerate}

\paragraph{Transition: from \gls{rwc} kinematics to the sense layer.}
The \gls{htlv} pipeline assumes that intruder state estimates are available to DAIDALUS at
the required rate and with modelled uncertainty. It does not quantify \emph{which} onboard
sensor provides the earliest cue, or how $P_d$ falls with range and lighting --- those are
properties of the perception stack. Section~\ref{sec:multisensor_benchmark} therefore presents
a synchronized Cosys-AirSim benchmark that isolates the ``sense'' path under identical
encounters, feeding the fusion and \gls{sum}/\gls{visa} rationale developed here.
